\chapter{John G. Thompson (2008)}

\section{Biographical Sketch}

John Thompson was born on 13 October 1932 in Ottawa, Kansas, USA. He studied at Yale University and the University of Chicago, where he completed his PhD in 1959 under the supervision of Saunders Mac Lane. He held positions at the University of Chicago, the University of Cambridge, and the University of Florida. Thompson received the Abel Prize (jointly with Tits) in 2008, the Fields Medal in 1970, the Wolf Prize in 1992, and the Steele Prize in 2000.

\section{High-Level View for a General Audience}

\subsection{What Did Thompson Do?}

Imagine you have a collection of shapes---triangles, squares, pentagons---and you want to understand all the ways you can rotate and reflect them so they land back on themselves. For a square, there are 8 such motions: 4 rotations and 4 reflections. For a regular pentagon, there are 10. This collection of motions is called a \emph{symmetry group}.

Now imagine the ultimate challenge: classify all possible symmetry groups, in any number of dimensions, with any number of elements. This is the problem of classifying \emph{finite simple groups}---the atomic building blocks of all finite symmetry groups. Thompson played a central role in this monumental achievement.

\paragraph{Thompson's Contribution:} Thompson proved that every finite group of odd order is solvable---the celebrated "odd order theorem" (also called the Feit--Thompson theorem). This was the first major step toward the classification of all finite simple groups. He also discovered several new sporadic simple groups, developed the theory of N-groups that became the template for the entire classification program, and uncovered deep connections between finite groups and modular forms that anticipated the monstrous moonshine phenomenon.

\subsection{Why Does It Matter?}

Symmetry groups are everywhere in mathematics and physics:

\begin{itemize}
\item They describe the structure of molecules in chemistry.
\item They are the foundation of particle physics (the Standard Model).
\item They are used in cryptography to secure communications.
\item They describe the symmetries of crystal structures.
\item They appear in art, architecture, and music.
\end{itemize}

The classification of finite simple groups is one of the greatest intellectual achievements of all time, comparable to the classification of chemical elements in the periodic table.

\section{The Origin of the Problem}

\subsection{Group Theory and Symmetry}

Group theory was born in the early nineteenth century through the work of Galois, Abel, Cauchy, and others. A \emph{group} is a set equipped with an associative binary operation, an identity element, and inverses. Groups capture the abstract notion of symmetry.

\begin{knowledgebridge}{What Is a Group?}
A \emph{group} is a set of actions that can be combined and reversed. For a square, the set of rotations and reflections forms a group: you can rotate by $90^\circ$ (an action), do it twice (combine actions), or undo it (take the inverse). Formally, a group requires three ingredients: an \emph{associative} operation (combining actions), an \emph{identity} element (doing nothing), and \emph{inverses} (undoing actions). Groups appear everywhere: the symmetries of a molecule, the set of moves on a Rubik's cube, and the fundamental laws of particle physics.
\end{knowledgebridge}

A group is \emph{simple} if it has no nontrivial normal subgroups---it cannot be decomposed into smaller pieces. Simple groups are the atoms of group theory: every finite group can be built from simple groups by repeated extensions (the Jordan--H\"older theorem).

\begin{knowledgebridge}{What Makes a Group Simple?}
A group is \emph{simple} if it cannot be broken into smaller, nontrivial pieces---like a chemical element that cannot be split by ordinary means. Just as every molecule is built from elements in the periodic table, every finite group can be built from simple groups (the Jordan--H\"older theorem). The classification of finite simple groups is the mathematical equivalent of listing all chemical elements.
\end{knowledgebridge}

\subsection{The Search for Simple Groups}

By the 1950s, mathematicians had identified several infinite families of finite simple groups:

\begin{itemize}
\item \textbf{Cyclic groups} $\mathbb{Z}_p$ of prime order.
\item \textbf{Alternating groups} $A_n$ for $n \geq 5$.
\item \textbf{Groups of Lie type}: families of matrix groups over finite fields, including $PSL(n, q)$, $PSp(2n, q)$, and the exceptional groups $G_2(q)$, $F_4(q)$, $E_6(q)$, $E_7(q)$, $E_8(q)$.
\end{itemize}

In addition, a handful of "sporadic" simple groups had been discovered, including the Mathieu groups $M_{11}$, $M_{12}$, $M_{22}$, $M_{23}$, $M_{24}$.

\subsection{The Classification Problem}

The central problem was: are there any more finite simple groups beyond these families? And can we prove that the list is complete?

This became known as the \emph{classification of finite simple groups} (CFSG), one of the most ambitious projects in the history of mathematics. The complete classification was announced in 1983 (with final corrections in 2004). It states that every finite simple group is:

\begin{enumerate}
\item A cyclic group of prime order;
\item An alternating group $A_n$ ($n \geq 5$);
\item A group of Lie type (including the Tits group);
\item One of 26 sporadic groups.
\end{enumerate}

\begin{analogy}{The Periodic Table of Symmetry}
Classifying finite simple groups is like creating a periodic table for symmetry. Just as Mendeleev discovered that chemical elements fall into families with predictable properties, mathematicians found that all finite simple groups belong to a few infinite families (cyclic, alternating, Lie type) plus 26 exceptional ``sporadic'' groups. The sporadic groups are like the noble gases---they do not quite fit the patterns of the families, yet they are essential for a complete picture.
\end{analogy}

\begin{figure}[htbp]
\centering
\begin{tikzpicture}[
    node distance=4mm and 4mm,
    box/.style={draw, thick, rounded corners, align=center, minimum width=2.8cm, minimum height=0.7cm, font=\small},
    titlebox/.style={draw, thick, rounded corners, align=center, font=\bfseries\small, minimum width=2.8cm},
    legend/.style={font=\footnotesize, align=center}
]

% Title
\node[font=\Large\bfseries, abelblue] at (6.5, 9.2) {Classification of Finite Simple Groups};

% --- Section 1: Cyclic ---
\node[box, fill=abelblue!15, draw=abelblue] (cyclic) at (2, 7.5) {$\mathbb{Z}_p$\\ {\footnotesize (prime order)}};
\node[legend, above=0pt of cyclic] {\textbf{Cyclic Groups}};

% --- Section 2: Alternating ---
\node[box, fill=abelred!15, draw=abelred, minimum width=3.5cm] (alt) at (6.5, 7.5) {$A_n$ ($n \ge 5$)\\ {\footnotesize alternating groups}};
\node[legend, above=0pt of alt] {\textbf{Alternating Groups}};

% --- Section 3: Lie Type ---
\node[titlebox, fill=abelgreen!15, draw=abelgreen, minimum width=8cm, minimum height=0.6cm] (lie) at (6.5, 5.8) {Groups of Lie Type};
\node[box, fill=abelgreen!8, draw=abelgreen, minimum width=7.5cm, minimum height=1.2cm] (classical) at (4.2, 4.5) {{\bf Classical:}\\ $PSL(n,q)$, $PSp(2n,q)$, $PSU(n,q)$, $PO(n,q)$};
\node[box, fill=abelgreen!8, draw=abelgreen, minimum width=7.5cm, minimum height=1.2cm] (exceptional) at (8.8, 4.5) {{\bf Exceptional:}\\ $G_2$, $F_4$, $E_6$, $E_7$, $E_8$};
\draw[abelgreen, thick] (lie.south) -- (classical.north);
\draw[abelgreen, thick] (lie.south) -- (exceptional.north);

% --- Section 4: Sporadic ---
\node[titlebox, fill=abelgray!15, draw=abelgray, minimum width=5cm] (spor) at (6.5, 2.5) {26 Sporadic Groups};
\node[box, fill=abelred!20, draw=abelred, minimum width=4cm, minimum height=0.8cm] (monster) at (6.5, 1.2) {{\LARGE\bfseries Monster}\\ $|M| \approx 8\times 10^{53}$};
\draw[abelgray, thick] (spor.south) -- (monster.north);

% Note about Tits group
\node[box, fill=yellow!15, draw=abelgray, minimum width=5cm] (tits) at (6.5, -0.3) {${}^2F_4(2)'$ (Tits group) is sporadic};
\draw[dashed, abelgray, thick] (monster.south) -- (tits.north);

% Connecting brackets on left
\draw[thick, abelblue] (0.2, 7.85) -- (0.2, 7.15);
\draw[thick, abelred] (0.2, 7.15) -- (0.2, 6.85);
\draw[thick, abelgreen] (0.2, 6.85) -- (0.2, 3.5);
\draw[thick, abelgray] (0.2, 3.5) -- (0.2, 0.4);

\end{tikzpicture}
\caption{The classification of finite simple groups: every finite simple group belongs to one of four families — cyclic groups of prime order, alternating groups of degree at least 5, groups of Lie type (classical or exceptional), or one of the 26 sporadic groups. The Tits group ${}^2F_4(2)'$, though closely related to Lie type, is classified as sporadic.}
\label{fig:finite_simple_groups_thompson}
\end{figure}

\section{Deep Insight into the Problem}

\subsection{Thompson's Odd Order Theorem}

Thompson's PhD thesis, written under the direction of Saunders Mac Lane, contained the germ of his most famous result. Working with Walter Feit, Thompson proved in 1963 that every finite group of odd order is solvable. This result, published as a 255-page paper in the Pacific Journal of Mathematics, is one of the longest and most important papers in the history of group theory.

\paragraph{The Feit--Thompson Theorem:} Every finite group of odd order is solvable. Equivalently, every non-abelian finite simple group has even order.

This theorem has profound consequences. It immediately implies that the smallest non-abelian simple group must have even order, and it launched the classification program by establishing that the structure of simple groups is tightly constrained.

\begin{bigidea}
The Feit--Thompson theorem says: every finite group of odd order is solvable. In plain terms, if a symmetry group has an odd number of elements, it cannot be a fundamental building block---it must be decomposable into simpler pieces. This was the crucial first step toward classifying all finite simple groups. The proof itself, spanning 255 pages, introduced revolutionary techniques (modular character theory, local analysis, signalizer functors) that became the template for the entire classification program.
\end{bigidea}

\subsection{Thompson's N-Group Theory}

Thompson developed the theory of \emph{N-groups}---finite groups in which every proper subgroup is solvable. He classified all simple N-groups in a monumental series of papers, showing that they are exactly the minimal simple groups (those whose proper subgroups are all solvable).

This work provided the foundation for the classification by reducing the general problem to the analysis of specific configurations.

\section{Biggest Contributions and New Tools}

\subsection{The Feit--Thompson Theorem}

The Feit--Thompson theorem (also called the Odd Order Theorem) is one of the deepest results in finite group theory. Its proof introduced several revolutionary techniques:

\begin{itemize}
\item \textbf{Character Theory}: The proof made extensive use of modular character theory, developed by Brauer. Thompson's "character-theoretic" approach to group theory became a standard method.

\item \textbf{Local Analysis}: The proof analyzed the structure of centralizers of involutions (elements of order 2), initiating the "local" approach to group classification.

\item \textbf{Signalizer Functors}: Thompson's concept of signalizer functors, which track the structure of subgroups commuting with elements of prime order, became a central tool in the classification.
\end{itemize}

\subsection{Thompson's Sporadic Groups}

Thompson discovered several sporadic simple groups:
\begin{itemize}
\item \textbf{Thompson's sporadic group Th}: A group of order 90,745,943,887,872,000, also called the \emph{Thompson group}.
\item \textbf{The Harada--Norton group HN}: Discovered by Thompson in collaboration with Harada and Norton.
\end{itemize}

\subsection{Thompson Series and Moonshine}

Thompson discovered connections between finite simple groups and modular forms that anticipated the monstrous moonshine phenomenon. The \emph{Thompson series} are modular functions associated to conjugacy classes of the Monster group.

\section{Impact of the Discovery in Science and Technology}

\subsection{Impact on Mathematics}

\begin{whyitmatters}
The classification of finite simple groups---to which Thompson contributed decisively---is one of the greatest intellectual achievements in mathematics, comparable to the classification of chemical elements or the mapping of the human genome. It provides a complete catalog of all possible finite symmetry patterns, with profound consequences for physics, cryptography, coding theory, and our understanding of the mathematical universe.
\end{whyitmatters}

Thompson's work transformed group theory and its connections:

\begin{itemize}
\item \textbf{Finite Group Theory}: The classification of finite simple groups, to which Thompson contributed the crucial first step and many essential techniques, is one of the greatest achievements of twentieth-century mathematics.

\item \textbf{Representation Theory}: Thompson's character-theoretic methods shaped modern representation theory.

\item \textbf{Number Theory}: Thompson's work on moonshine forged deep connections between finite groups and modular forms, opening an entirely new area of research.
\end{itemize}

\subsection{Impact on Physics}

\begin{itemize}
\item \textbf{Particle Physics}: The classification of simple groups explains the symmetry groups of fundamental interactions. The Standard Model's gauge group $SU(3) \times SU(2) \times U(1)$ is a product of simple groups and abelian factors.

\item \textbf{String Theory}: Exceptional groups and the Monster group, related to Thompson's work, appear in string theory.

\item \textbf{Condensed Matter Physics}: Symmetry groups are used to classify phases of matter, and the classification of simple groups provides the complete list of possible symmetry types.
\end{itemize}

\subsection{Impact on Technology}

\begin{itemize}
\item \textbf{Cryptography}: Group theory is used in cryptography, including the use of group structures in elliptic curve\index{Elliptic curves} cryptography and lattice-based cryptography.

\item \textbf{Coding Theory}: Group-theoretic codes, including group codes and Reed--Muller codes, use the classification of simple groups.

\item \textbf{Quantum Computing}: The theory of quantum error-correcting codes uses group theory, and the classification of simple groups provides constraints on possible code structures.
\end{itemize}

\section{Current Developments}

\subsection{The Second-Generation Classification}

Work continues on simplifying and verifying the classification of finite simple groups:

\begin{itemize}
\item \textbf{The GLS Project}: Gorenstein, Lyons, and Solomon are writing a series of volumes providing a complete and unified proof of the classification.

\item \textbf{Aschbacher's Contributions}: Aschbacher has made major contributions to simplifying portions of the classification proof.

\item \textbf{Computer Verification}: Formal verification of parts of the classification using proof assistants is an ongoing project.
\end{itemize}

\subsection{Monstrous Moonshine}

The connection between finite simple groups and modular forms that Thompson pioneered continues to be an active area:

\begin{itemize}
\item \textbf{Moonshine Beyond the Monster}: Extensions of moonshine to other groups, including the Mathieu group $M_{24}$ and the Conway group, have been discovered.

\item \textbf{Physical Moonshine}: Connections between moonshine and string theory, including the K3 sigma model, are an active research area.
\end{itemize}

\section{Conclusion}

John Thompson transformed our understanding of symmetry. His work on the odd order theorem, the theory of N-groups, and the classification of simple groups provided the conceptual framework for one of the greatest intellectual achievements in the history of mathematics: the complete classification of finite simple groups.

The Abel Prize citation recognized Thompson (jointly with Tits) "for their profound achievements in algebra and in particular for shaping modern group theory." Thompson's work has shaped not just algebra, but geometry, representation theory, number theory, and mathematical physics.

\section{Behind the Breakthrough: The Human Story}
Thompson's Feit-Thompson proof (establishing the Odd Order Theorem) was 255 pages long, occupying an entire issue of the Pacific Journal of Mathematics. At the time, this was completely unprecedented, breaking the convention of short papers and paving the way for the classification of finite simple groups.

\section{Pedagogical Metaphors and Lineage}
\subsection{Signature Metaphor}
``Classifying the fundamental building blocks of all algebraic symmetry by proving that odd-order groups must be decomposable.''

\subsection{Mathematical Lineage}
Thompson was advised by Saunders Mac Lane (co-founder of category theory). Thompson's doctoral students include R. L. Griess (who constructed the Monster group) and many other leading group theorists.

\section{Applied Science and Computational Algorithms}
\subsection{Key Takeaways for Applied Science}
Group theory is used in crystallography (analyzing atomic configurations), structural chemistry (molecular symmetries), and in coding theory for designing algebraic error-correcting codes.

\subsection{Algorithms and Numerical Implementations}
Coset enumeration algorithms (such as the Todd-Coxeter algorithm) and the Schreier-Sims algorithm for permutation groups (implemented in computer algebra systems like GAP and Magma) are used to compute simple group representations.

\section{The Research Frontier}
Completing the second-generation proof of the Classification of Finite Simple Groups, and further exploring the connections between finite groups and modular forms in moonshine theory.

\section{Guided Exercises and Challenges}
\begin{enumerate}[label=\textbf{\arabic*.}]
  \item Explain the significance of the Feit-Thompson Odd Order Theorem in the classification of finite simple groups.
  \item What are N-groups and why were they important for the classification program?
  \item Describe Thompson's contributions to the discovery of sporadic simple groups.
\end{enumerate}

\section{Curriculum Map and Further Reading}
For students wishing to delve deeper into the mathematical machinery underlying these breakthroughs, the following learning path is recommended:
\begin{itemize}
  \item Michael Aschbacher, \emph{The Finite Simple Groups}, Cambridge University Press.
  \item D. Gorenstein, \emph{Finite Simple Groups: An Introduction to Their Classification}, Plenum.
\end{itemize}

\section*{Bibliography}

\begin{enumerate}
\item W. Feit and J. G. Thompson, "Solvability of groups of odd order," Pacific Journal of Mathematics 13 (1963), 775--1029.
\item J. G. Thompson, "A simple group of order 90,745,943,887,872,000," (with R. M. Guralnick), in preparation.
\item J. G. Thompson, "Finite groups with fixed-point-free automorphisms of prime order," Proceedings of the National Academy of Sciences 45 (1959), 578--581.
\item J. G. Thompson, "Nonsolvable finite groups all of whose local subgroups are solvable," Bulletin of the American Mathematical Society 74 (1968), 383--437.
\item J. G. Thompson, "Some finite groups which appear as Galois groups of algebraic number fields I--III," 1972--1984.
\item D. Gorenstein, "Finite Simple Groups: An Introduction to Their Classification," Plenum, 1982.
\item M. Aschbacher, "The Classification of Finite Simple Groups," American Mathematical Society, 2004.
\end{enumerate}
